\documentclass[11pt,a4paper]{article}
\usepackage[margin=24mm]{geometry}
\usepackage{amsmath,amssymb,amsthm,booktabs,microtype}
\setlength{\emergencystretch}{3em}
\newtheorem{theorem}{Theorem}
\newtheorem{lemma}[theorem]{Lemma}
\newcommand{\GL}{\operatorname{GL}}
\newcommand{\SL}{\operatorname{SL}}
\newcommand{\diag}{\operatorname{diag}}
\newcommand{\F}{\mathbb F}
\title{Involutions and intersections in groups with determinant $\pm1$}
\author{Human Compression project: N012 recovery}
\date{A component of the thesis recovery; distinct review pending}
\begin{document}\maketitle

Let $q=p^f$ be odd and put
\[
 H_n(q)=\{g\in\GL_n(q):\det g\in\{1,-1\}\}.
\]
We recover the constructive and polytope assertions of Markovskaya's
Theorems 4.2.1 and 4.2.2, and the part for dimensions at least five of the
quotient assertions in Chapter 4 of [T]. The matrices and most root words
are source constructions. Their generation arguments are replaced by
the complete root-parameter theory [U]. A single five-dimensional
calculation supplies the intersection property for every construction
and every scalar quotient. The six-dimensional construction for $q>3$
is a normalization of the already recovered journal matrices [L].

We use $t_{ij}(a)=I+aE_{ij}$, $x^y=yxy^{-1}$ and
$[x,y]=xyx^{-1}y^{-1}$. Write $\tau e_i=e_{n+1-i}$ and
$\mu e_i=e_{i+1}$, where $e_{n+1}=e_1$. For $\epsilon\in\{1,-1\}$,
let $C_\epsilon$ reverse the first $n-1$ coordinates and send
$e_n$ to $\epsilon e_n$. Then $\tau C_\epsilon$ is the monomial
cycle with $e_i\mapsto e_{i+1}$ for $i<n$ and
$e_n\mapsto\epsilon e_1$.

\section{The common generation step}
The following complete statements, with their proofs, are acquired
from [U]. A monomial cycle and one nonidentity root on two cyclically
consecutive coordinates reach every root subgroup: conjugate the seed
around the cycle and apply
$[t_{ij}(a),t_{jk}(b)]=t_{ik}(ab)$ along simple directed paths.
If a subgroup $M\le\GL_n(q)$, $n\ge3$, meets every root subgroup,
its root-generated subgroup is diagonally conjugate to $\SL_n(K)$
for a subfield $K\le\F_q$. In particular it contains all
determinant-one diagonal sign matrices. Opposite-root parameters
$a,b$ have product $ab\in K$. Finally a primitive $u\in\F_q^*$
satisfies $\F_p(u^2)=\F_q$. These are the particular forms needed
below; none is inferred from testing finitely many fields.

Let $A=A(n,a)$ have off-diagonal entries
$A_{21}=A_{n-1,n}=a\ne0$ and no others. For $n=5,6$ its diagonal
is negative at $1,n$ and positive elsewhere. For $n\ge7$ it is
negative at $2,3,n-2,n-1$ and positive elsewhere.
In all cases $A^2=I$ and $A\tau=\tau A$.
Put $D=\diag(1,-1,-1,1,\ldots,1)$. Whenever all roots have
already been reached, $D\in M$ and direct multiplication gives
\[
 [D,A]=\begin{cases}t_{21}(2a),&n=5,6,\\
                     t_{21}(-2a),&n\ge7.
        \end{cases}                                      \tag{1}
\]
This commutator isolates a root only after root recognition has
supplied $D$; it is not used to justify the initial seed.

For $B=\tau$ or $-\tau$ and $E=BC_\epsilon$, transporting
$t_{21}(b)$ first by $B$ and then by $E^2$ gives $t_{12}(b)$.
The two wraparound weights in the latter transport cancel.
Consequently (1), with primitive $a=u$, forces
$4u^2\in K$ and hence $K=\F_q$. It follows that
$\SL_n(q)\le M$. If one generating matrix has determinant $-1$
and all have determinant in $\{1,-1\}$, then $M=H_n(q)$.

\section{Five coordinates}
For $q>3$ choose a primitive $u\ne2$. There are at least two
primitive elements in the cyclic group of order $q-1\ge4$, so
this choice is possible. For $q=3$ take $u=1$; in that prime field
the conclusion about the parameter field is automatic. Use
\[
 A=A(5,u),\qquad B=-\tau,\qquad C=C_1,\qquad E=BC=-\mu.
\]
These are involutions, $A$ and $B$ commute, and $\det B=-1$.
The source words on pages 58--59 give a nonzero adjacent root as
follows. Define
\[
 X=[A,A^E],\quad D_0=[A,X^E],\quad
 Y=[X^E,D_0^{E^{-1}}],\quad T=Y^{E^3},\quad J=[A^{E^2},T].
\]
Using $E_{ij}E_{ab}=\delta_{ja}E_{ib}$, their complete nonidentity
entries are
\begin{align*}
 X&=t_{31}(-u^2)t_{41}(u^2),\\
 D_0&=t_{41}(u^4-u^3)t_{42}(u^3)t_{51}(-u^3)t_{52}(-2u^2),\\
 D_0^{E^{-1}}&=t_{31}(u^3)t_{41}(-2u^2)
                  t_{35}(u^4-u^3)t_{45}(-u^3),\\
 Y&=t_{32}(-u^6+u^5)t_{42}(u^5),\\
 T&=t_{15}(-u^6+u^5)t_{25}(u^5),\\
 J&=t_{15}(u^6)t_{25}(-2u^5),\\
 T^2J&=t_{15}\bigl(u^5(2-u)\bigr).
\end{align*}
The factors in each displayed product commute. The last parameter
is nonzero by the choice of $u$, including $u=1$ in $\F_3$.
The cycle $E$ propagates this root to all roots. The common
generation step gives $\langle A,B,C\rangle=H_5(q)$.

\section{Seven coordinates and the uniform construction}
For $n\ge7$ take primitive $u\in\F_q^*$ and set
\[
 A=A(n,u),\quad B=\tau,\quad C=C_\epsilon,\qquad
 \epsilon=\begin{cases}1,&n\equiv0,3\pmod4,\\
                       -1,&n\equiv1,2\pmod4.
          \end{cases}
\]
Then $\det C=-1$, and $A,B,C$ are involutions with $AB=BA$.
Put $E=BC$ and
\[
 X=[A,A^E],\quad Y=(XA^{E^2})^2,\quad T=Y^E,
 \quad Z=[T,X],\quad R=[A,Z],\quad V=R^B,\quad W=[Z,V].
\]
For every $n\ge8$, direct multiplication gives
\begin{align*}
 X&=t_{31}(-u^2)t_{n-1,1}(\epsilon u^2),\\
 Y&=t_{32}(-u^3)t_{41}(u^3)t_{42}(-u^4)t_{n-1,2}(\epsilon u^3),\\
 T&=t_{43}(-u^3)t_{52}(u^3)t_{53}(-u^4)t_{n3}(\epsilon u^3),\\
 Z&=t_{41}(u^5)t_{51}(u^6)t_{n1}(-\epsilon u^5),\\
 R&=t_{n-1,1}(-\epsilon u^6),\qquad
 V=t_{2n}(-\epsilon u^6),\qquad W=t_{21}(-u^{11}).       \tag{2}
\end{align*}
Here is a uniform verification of the range of these identities.
For the first two products use the eight distinct labels
$1,2,3,4,5,n-2,n-1,n$. The diagonal negatives of $A,A^E,A^{E^2}$
are respectively
\[
 \{2,3,n-2,n-1\},\quad\{3,4,n-1,n\},\quad\{4,5,n,1\}.
\]
Their respective off-diagonal pairs are
\[
 (21:u,\ n-1,n:u),\quad(32:u,\ n,1:\epsilon u),
 \quad(43:u,\ 12:u).
\]
All remaining coordinates are fixed. Multiplication on these labels
gives $X,Y$. Conjugating the four entries of $Y$ shifts each index
one step without wraparound and gives $T$. The matrix-unit rule then
gives $Z,R$; reversal transports $R$ to $V$, and a final commutator
gives $W$. Both possible values of $\epsilon$ satisfy the same
identities. The labels are distinct also when $n=8$, so no unexamined
small boundary is hidden in (2).

For $n=7$ we have $\epsilon=1$. The same definitions give the same
$X,T,Z$, but the overlap $n-2=5$ changes the remaining products to
\begin{align*}
 R&=t_{51}(-2u^6)t_{61}(-u^6),\\
 V&=t_{27}(-u^6)t_{37}(-2u^6),\\
 W&=t_{21}(-u^{11})t_{31}(-2u^{11}),\\
 Q:=[W,W^E]&=t_{31}(-u^{22})t_{41}(-2u^{22}),\\
 J:=[A,Q]&=t_{31}(2u^{22}),\qquad JQ^2=t_{41}(-4u^{22}),\\
 [(JQ^2)^{E^2},V]&=t_{67}(8u^{28}).                    \tag{3}
\end{align*}
These are identities over $\mathbb Z[u]$, so the last parameter is
nonzero in every odd characteristic. The square in $JQ^2$ is necessary:
the printed cancellation on source page 67 omits it. We correct that
individual expression, rather than infer a failure of the theorem.
The adjacent root in (2) or (3), followed by the common generation
step, proves $\langle A,B,C\rangle=H_n(q)$ for all $n\ge7$.

\section{Six coordinates}
First let $q>3$. We choose $v\in\F_q^*$ such that
\[
 -v\text{ is a square},\qquad \F_p(v)=\F_q,
 \qquad h=v^2+4v\ne0.                                  \tag{4}
\]
If $q\equiv3\pmod4$, choose primitive $v\ne-4$.
There are at least two primitive choices; both $v$ and $-1$ are
nonsquares, so $-v$ is square. If $q\equiv1\pmod4$, choose
$v=u^2\ne-4$ with primitive $u$. The primitive-square result gives
the field assertion. At $q=5$ the choice is $v=4$, whereas $-4=1$.
For $q>5$, the cyclic order $q-1\ge8$ has at least four primitive
elements, giving at least two distinct squares, so one avoids $-4$.
In this case $-1$ is square as well. This proves (4).

Take $A=A(6,-1)$, $C=C_1$, and
\[
 B_0=\diag(-v^{-1},-v^{-1},-v^{-1},1,1,1)\tau,
 \qquad B=\lambda B_0,\quad\lambda^2=-v.
\]
Then $B_0^2=-v^{-1}I$, so $B^2=I$; also $AB=BA$.
Their determinants are $\det A=\det C=1$ and $\det B=-1$.
Conjugation by $B$ equals that by $B_0$, and conjugation by
$E=BC$ equals that by $E_0=B_0C$.

The complete six-dimensional journal word from [L] is a
Laurent-polynomial identity in $v$; its derivation does not require
$v$ primitive. In the present notation it reads
\begin{align*}
 X=[A,A^E]&=t_{31}(-1)t_{51}(-v),\\
 Q=[A^{E^4},X]&=t_{51}(2v)t_{61}(-v^2),\\
 [A,Q]&=t_{51}(-h),\\
 (A^{E^4}t_{51}(h))^2&=t_{61}(vh),\qquad
 t_{61}(vh)^E=t_{12}(vh).
\end{align*}
Every element on the left is in the generated group, since
$t_{51}(h)=[A,Q]^{-1}$. By (4) the last root is nontrivial.
Propagation gives every root. Now (1) gives $t_{21}(-2)$, and
\[
 \bigl(t_{21}(-2)^B\bigr)^{E^2}=t_{12}(2v).
\]
Thus the parameter field contains $-4v$ and hence $v$; (4) makes
it the full field. The determinant calculation finishes generation
of $H_6(q)$.

For $q=3$ use the original finite-field marking
\[
 A=A(6,1),\qquad B=\tau,\qquad C=C_{-1}.
\]
It suffices to give one adjacent elementary matrix. The following
literal word in $a=A,b=B,c=C$, read left to right, has 131 letters
and evaluates over $\F_3$ to $t_{21}(1)$; the line breaks have no
algebraic meaning:
\begin{center}\small\ttfamily
abcbcacbcacabcacacacbcbcacabcbcac\par
bacbcacacacabcacabcacbcacbcacbcac\par
babcacbcacabcacacacbcacbcacbcacbc\par
acbcacbcacabcacacacbcacbcacbcacb
\end{center}
The exact word is also retained as an integer list. Its verification
consists solely of multiplying the three displayed $6\times6$ matrices
modulo three; the accompanying standard-library checker performs these
131 products and compares every entry with $I+E_{21}$. No group-order
calculation or completeness of the search that found the word is a
premise. Conjugating around $BC$ and taking elementary commutators
gives every root over the prime field. Hence the group contains
$\SL_6(3)$; since $\det B=-1$, it is $H_6(3)=\GL_6(3)$.

\section{One intersection calculation}
The preceding constructions share a five-dimensional invariant
subspace for $A,C$. Put $a_0=A_{21}=A_{n-1,n}\ne0$ and let
$\sigma=-1$ for $n=5,6$, $\sigma=1$ for $n\ge7$.
Write $\delta=-\sigma$ and $\epsilon=C_{nn}$. On the ordered basis
\[
 x=e_1,\quad y=e_2,\quad z=e_{n-2},\quad w=e_{n-1},\quad t=e_n,
\]
the matrix $F=AC$ acts by
\[
 x\mapsto\delta w,\quad y\mapsto\delta z,\quad
 z\mapsto\delta y,\quad w\mapsto\sigma x+a_0y,
 \quad t\mapsto\epsilon\sigma t+\epsilon a_0w.           \tag{5}
\]
On the complementary coordinates it is a signed permutation whose
permutation cycles have length at most two, so its fourth power is
identity. Within (5), squaring gives
\begin{align*}
 F^2x&=-x+\delta a_0y,& F^2w&=-w+\delta a_0z,\\
 F^2y&=y,&F^2z&=z,\\
 F^2t&=t+\sigma a_0w+\epsilon\sigma a_0x+\epsilon a_0^2y.
\end{align*}
Squaring again therefore proves, for every dimension in question,
\[
 F^4=I+N,\qquad
 N=a_0^2(\epsilon E_{2n}-E_{n-2,n}),\quad N^2=0.         \tag{6}
\]
In particular $N\ne0$. We next prove that a power $F^j$ is monomial
if and only if $4p$ divides $j$, in which case it is identity.
For $j=4l$, (6) gives $F^j=I+lN$. Its last column has three nonzero
entries when $l\ne0$ in $\F_p$, and it is identity otherwise.
The other residues have two fixed nonzero entries that cannot be
altered by right multiplication by $I+lN$, which changes only column $n$:
\[
\begin{array}{c|c|c}
\text{matrix}&\text{row}&\text{two entries (column:value)}\\\hline
F&2&n-2:\delta,\quad n-1:a_0\\
F^2&2&1:\delta a_0,\quad2:1\\
F^{-1}=CA&n-2&1:a_0,\quad2:\delta.
\end{array}
\]
The last row covers $F^{4l+3}=F^{-1}(I+(l+1)N)$.
This proves the assertion for all integers $j$, as well as
$|AC|=4p$.

Since $A,C$ are involutions, every element of $\langle A,C\rangle$
is $F^j$ or $F^jC$. The matrix $C$ is monomial, so the only monomial
elements in that dihedral group are $I,C$. Both $B$ and $C$ are
monomial in every construction. Hence
\[
                  \langle A,C\rangle\cap\langle B,C\rangle
                         =\langle C\rangle.             \tag{7}
\]
The three involutions are distinct: $A$ is not monomial, while $B,C$
have different permutation parts. With the order $(A,C,B)$, their
end involutions commute. The standard rank-three group--polytope
correspondence applied to (7) therefore gives a regular polytope.

The order of $BC$ is also explicit. For $n=5$ it is $2n$ since
$BC=-\mu$. For $n\ge7$ its $n$th power is $\epsilon I$, and
no smaller positive power has identity permutation part. Thus its
order is $n$ or $2n$. In the six-dimensional normalized construction,
the product of the six monomial weights is one: its determinant is
$-1$ and the permutation is a six-cycle of sign $-1$. Its order is
therefore six. In the exceptional $q=3$ construction, $(BC)^6=-I$,
so its order is twelve. In every case the type is $[4p,n]$ or
$[4p,2n]$.

\begin{theorem}
For odd $q$ and $n\ge2$, $H_n(q)$ admits three generating involutions
with two commuting if and only if $n\ge5$. For every $n\ge5$ it is
the automorphism group of a regular polytope of rank three and type
$[4p,n]$ or $[4p,2n]$.
\end{theorem}
\begin{proof}
The positive constructions and (7) prove every assertion for $n\ge5$.
For $n=2,3,4$, the complete symmetric-square obstruction proved in
the acquired central-quotient component [Z] excludes even the required
generation. Its hypotheses are exactly odd $q$ and the determinant-sign
group. We retain that proof rather than duplicate it.
\end{proof}

\begin{theorem}
Let $n\ge5$ and let $S$ be any subgroup of scalar matrices in $H_n(q)$.
The images of the displayed marking give $H_n(q)/S$ a regular polytope
of type $[4p,n]$ or $[4p,2n]$. For the full scalar subgroup
$S=Z(H_n(q))$, the type is $[4p,n]$.
\end{theorem}
\begin{proof}
Suppose $x\in\langle A,C\rangle$ and $y\in\langle B,C\rangle$
have the same image. Then $x=ys$ for a scalar $s\in S$, so $x$ is
monomial and therefore $x\in\{I,C\}$. This proves the intersection
property in the quotient. Moreover $\langle A,C\rangle\cap S=1$,
since $C$ is not scalar. Its order $4p$ and the distinctness of the
three involutions persist. The permutation part of $BC$ is an
$n$-cycle, so its image has order divisible by $n$. Its $n$th power
is $I$ or $-I$, hence the quotient order is $n$ or $2n$; in the full
scalar quotient it is exactly $n$. Generation passes to the quotient.
\end{proof}

\section{Attribution and verification boundary}
The recovered $H_n(q)$ classification and its polytope types are
Markovskaya's results [T, pp.~56--69]. The matrices and seed words in dimensions five and at least seven
are hers. Formula
(3) repairs a missing square on the actual printed page 67.
The six-dimensional normalized marking comes from [L], with a
different parameter choice and scalar normalization; the prime-field
word supplies the remaining case. The root theory is the acquired
classical theory proved in [U]. The common five-dimensional argument
expands the source's monomial-intersection calculation and makes its
descent to every scalar quotient explicit. We make no literature-wide
novelty claim for these arguments.

The exact polynomial checker verifies the displayed finite words,
the uniform eight-label computation and the four sign choices in
(5). The written matrix-unit proof, not a finite sample of dimensions,
justifies the uniform ranges. The 131-letter prime-field identity has
its separate complete multiplication certificate. Neither is a Lean
formalization. This component does not settle the remaining small
quotient cases or the full Chapter 3 routing audit, and gives no
credit for the whole thesis, an open question or a measured compression factor.

\begin{thebibliography}{9}\small
\bibitem{T}[T] I.~A.~Markovskaya, \emph{Linear groups over finite fields
and over the ring of integers as automorphism groups of regular
polytopes of rank 3}, thesis, Siberian Federal University, 2025,
80-page supplied edition; supervisor Ya.~N.~Nuzhin.
\bibitem{L}[L] \emph{Three involutions in finite general linear groups},
reviewed recovery of I.~A.~Markovskaya and Ya.~N.~Nuzhin,
\emph{On generation of the groups $\GL_n(q)$ and $\operatorname{PGL}_n(q)$
by three involutions, two of which commute}, Trudy Inst. Mat. Mekh.
UrO RAN 31(4) (2025), 247--259.
\bibitem{U}[U] \emph{Root parameters and fixed-space obstructions},
complete reviewed common paper accompanying [L].
\bibitem{Z}[Z] \emph{Central quotients of the determinant-sign subgroup},
reviewed N012 component, including its complete low-dimensional
generation obstruction and exact centre calculation.
\end{thebibliography}
\end{document}
